\documentclass[11pt]{article}
\usepackage[margin=1.1in]{geometry}
\usepackage{amsmath,amssymb,amsthm}
\usepackage{microtype}

\theoremstyle{plain}
\newtheorem{theorem}{Theorem}
\newtheorem{proposition}[theorem]{Proposition}
\theoremstyle{definition}
\newtheorem{example}[theorem]{Example}
\newtheorem{remark}[theorem]{Remark}

\newcommand{\ZZ}{\mathbb Z}
\newcommand{\vlam}{v_\lambda}

\title{Confirmation on the agreement locus\\
{\large three worked examples of a mechanical failure mode}}
\author{W.~Dahn}
\date{}

\begin{document}
\maketitle

\begin{abstract}
\noindent A conjectured formula is tested, confirmed repeatedly, and believed.
It is false. Every test happened to lie in the set where the conjecture agrees
with the truth --- and that set was not chosen adversarially but by
computational convenience, because the cases that were cheap to compute were
exactly the cases where the two formulas coincide. We give three examples from
a single research programme, all found within one week of each other, and
isolate the structural reason: \emph{the constraint that makes a case
affordable is often the same constraint that makes the conjecture true there}.
The third example is the sharpest, because the false statement was not a fit
but a \emph{theorem} --- correctly proved from a premise that had itself been
derived only on the agreement locus. We propose a cheap test.
\end{abstract}

\section{The failure mode}

Let $C$ be a conjecture, $T$ the truth, and $A=\{x: C(x)=T(x)\}$ their agreement
locus. Confirming $C$ at points of $A$ produces no evidence for $C$: the
observations are equally consistent with $T$. This is obvious stated abstractly
and invisible in practice, because $A$ is unknown --- one does not learn what
$A$ is until $C$ has already failed somewhere.

What makes the failure mode mechanical rather than merely possible is that test
sets are rarely chosen at random. They are chosen where computation is
feasible. If feasibility and membership in $A$ are governed by the same
condition, then every affordable test is uninformative, and no amount of
confirmation helps. That is what happened three times below.

\section{Example one: a factorial that was not there}

The setting is a family of matrices $D$ over $\ZZ[\zeta_p]$ indexed by a finite
set of ``sections'', with $\lambda=1-\zeta_p$ and $\vlam$ the corresponding
valuation. The conjecture, fitted to thirty-eight measured points, was
\begin{equation}\label{eq:C}
\vlam\bigl(\operatorname{tr}D^m\bigr)
=\vlam(q)+m+[m\text{ odd}]+\vlam(m!),
\end{equation}
the minimum being over sections. It survived fourteen rounds of testing,
including a nine-step check against Legendre's formula for $v_p(m!)$ that
matched at all twenty-six exponents of one case.

At $q=3$ the section space has exactly $81$ elements --- small enough to
enumerate --- and the minimum is therefore decidable. Enumerating it gives
\[
4,\;8,\;8,\;\mathbf{12},\;12,\;\mathbf{16},\;\mathbf{16},\;20,\;20,\;
\mathbf{24},\;24 \qquad(m=2,\dots,12),
\]
against \eqref{eq:C}'s $4,8,8,10,12,14,14,20,20,22,24$. The truth exceeds the
conjecture by $2$ at $m=5,7,8,11$, and fits
\[
T(m)=2\bigl(m+[m\text{ odd}]\bigr)
\]
at all nineteen exponents up to $m=20$, where \eqref{eq:C} fails at eleven.

Now the agreement locus. Since $2v_3(m!)=m-s_3(m)$ with $s_3$ the base-$3$
digit sum, the two expressions differ by $s_3(m)+[m\text{ odd}]-2$, so
\[
A=\{\,m:\ s_p(m)+[m\text{ odd}]=2\,\}.
\]
Every prime power $m=p^{j}$ has $s_p(m)=1$ and is odd, hence lies in $A$. So
does $m=q$ for every prime $q=p$. \emph{Every test that had been run was at
$m=q$, at $m=p^{j}$, or at two composite ring sizes} --- all inside $A$.

Why those? Because the analytic machinery available decomposes
$\operatorname{tr}D^m$ into classes indexed by the divisors of $m$, and at
prime powers that decomposition collapses to a single class. Prime powers were
the cases that could be computed at all. Feasibility and $A$ coincided.

\section{Example two: a completeness claim}

The same programme proved a ``sieve theorem'' giving $|T|$ linear relations
among $\tau(m)$ classes, and then asked whether those were all of them. The
test was a rank computation: measure the classes across many sections and read
off the nullity. Reported: nullity $=|T|$ in ten cells, hence the sieve is the
whole relation space.

The cells were $(p,m)=(3,3),(3,6),(3,9),(3,15),(3,27),(3,81),(5,5),(5,25),
(7,7),(7,49)$. Every one satisfies $p\mid m$. Choosing instead cells with
$|T|=0$ --- where the claim predicts \emph{no} relations, so a stray one cannot
hide --- gives nullity $1$ at $(3,2)$, $(3,4)$, $(3,10)$, $(5,4)$ and $(5,6)$
alike. The correct count is
\[
\text{nullity}=|T|+[\,p\nmid m\,].
\]
The extra relation is elementary once seen: one of the classes is \emph{empty}
when $p\nmid m$, so its sum vanishes vacuously, and the sieve --- which counts
relations forced by cancellation --- does not count it.

Again the locus was not chosen adversarially. The computational shortcut that
made the classes affordable requires a divisibility condition equivalent to
$p\mid m$. The affordable cells were exactly the cells where the missing
relation cannot occur.

\section{Example three: a theorem inheriting the blind spot}

The third instance is the one worth the paper. From the (incorrect) count
$\mathrm{free}(m)=\tau(m)-|T|$, the programme \emph{proved} a characterisation:
$\mathrm{free}(m)=1$ if and only if $m$ is a power of $p$. The proof is a
correct divisor computation. But with the corrected count
$\mathrm{free}(m)=\tau(m)-|T|-[\,p\nmid m\,]$, every prime $\ell\neq p$ has
$\tau=2$, $|T|=0$ and correction $1$, hence $\mathrm{free}=1$ --- verified at
$(3,5)$, $(3,7)$, $(5,3)$, $(5,7)$, $(7,3)$, none of which is a power of $p$.

\begin{remark}
The algebra was sound; the premise was not. A proof inherits the blind spot of
whatever it starts from, and a formula derived only on an agreement locus
carries that locus into every theorem built on it. Formalisation would not have
caught this: the Lean modules in that programme state their arithmetic inputs
as hypotheses, and the hypothesis here was false.
\end{remark}

\section{A cheap test}

For a conjecture $C$ predicting a quantity, run it where $C$ predicts something
\emph{degenerate} --- zero relations, an empty set, a vanishing invariant.
Degenerate predictions have no room to absorb an error, so a discrepancy shows
immediately. In example two the entire refutation was five cells with $|T|=0$.

More generally, before trusting a confirmation, ask the two questions:

\begin{enumerate}
\item \textbf{Was the test set chosen independently of the conjecture?} If the
cases were selected because they were computable, identify the condition that
made them computable and check whether it appears in $C$.
\item \textbf{Can the conjecture be made to predict something degenerate?} If
so, test there first; those cases are usually cheap for the same reason they
are informative.
\end{enumerate}

A useful negative control: in the same programme, several results were
\emph{not} affected --- exact identities checked against honest enumeration,
and a determinant congruence proved unconditionally and verified by exhaustion
at $q=3$. The distinguishing feature is not rigour but provenance: a fitted
formula tested on convenient cases is exposed; an identity verified against a
direct computation is not.

\section*{Provenance}

All three examples arise in the author's work on Heisenberg--Weyl determinant
congruences; the underlying computations and their certificates are archived
with that work, together with the witness programs that produced both the
original confirmations and the refutations. The refutations were found in
passes 519, 522 and 523 of that programme, within a single week, after the
failure mode was named in pass 521.

\end{document}
